\begin{tikzpicture}[font=\small\sffamily,text=ink,
 card/.style={draw=ink!20,rounded corners=4pt,fill=wash,minimum width=7.65cm,minimum height=4.8cm},
 block/.style={draw=teal!45,fill=white,rounded corners=2pt,align=center,text width=1.95cm,minimum height=.75cm,font=\footnotesize\sffamily},
 flow/.style={-{Latex[length=1.7mm]},draw=ink!60,line width=.7pt}]
\foreach \x/\y in {-4/0,4/0,-4/-5.45,4/-5.45} {\node[card] at (\x,\y) {};}
\node[anchor=north west,text width=7cm,align=left,font=\small\sffamily\bfseries] at (-7.55,2.2) {DIGITAL PDF OR SCAN\\ENOUGH EMBEDDED TEXT};
\pic[scale=.8] at (-7.3,1.05) {dqtext};
\node[align=left,text width=5cm,font=\footnotesize\sffamily] at (-2.55,.4) {No specialist selected.\\Enough embedded text selects the text route, even if the PDF looks scanned.};
\node[block] (a1) at (-6.5,-1) {Primary\\reader};
\node[block] (a2) at (-4,-1) {Primary\\checks};
\node[block] (a3) at (-1.5,-1) {Return its\\JSON};
\draw[flow] (a1)--(a2);\draw[flow] (a2)--(a3);
\node[align=left,text width=7cm,font=\footnotesize\sffamily,text=ink!70] at (-4,-1.9) {Page images may accompany text. No general quorum is needed for this single reading.};
\node[anchor=north west,text width=7cm,align=left,font=\small\sffamily\bfseries] at (.45,2.2) {SHORT SCAN\\PARSER UNAVAILABLE};
\pic[scale=.8] at (.7,1.05) {dqscan};
\node[align=left,text width=5cm,font=\footnotesize\sffamily] at (5.45,.4) {No sufficient text; $n\leq12$.\\No parse candidate, but primary and one-shot may still succeed.};
\node[block] (b1) at (1.5,-1) {Primary +\\one-shot};
\node[block] (b2) at (4,-1) {Check each\\separately};
\node[block] (b3) at (6.5,-1) {Quorum if\\both succeed};
\draw[flow] (b1)--(b2);\draw[flow] (b2)--(b3);
\node[align=left,text width=7cm,font=\footnotesize\sffamily,text=ink!70] at (4,-1.9) {One successful reading returns directly. No successful reading raises an error.};
\node[anchor=north west,text width=7cm,align=left,font=\small\sffamily\bfseries] at (-7.55,-3.25) {LONG SCAN\\SUCCESSFUL CACHED PARSE};
\pic[scale=.8] at (-7.3,-4.4) {dqscan};
\node[align=left,text width=5cm,font=\footnotesize\sffamily] at (-2.55,-5.05) {No sufficient text; e.g. $n=20$.\\Try primary, then the parse reading. A successful parse takes priority.};
\node[block] (c1) at (-6.5,-6.45) {Read primary,\\then parse};
\node[block] (c2) at (-4,-6.45) {Check each\\separately};
\node[block] (c3) at (-1.5,-6.45) {Return\\checked parse};
\draw[flow] (c1)--(c2);\draw[flow] (c2)--(c3);
\node[align=left,text width=7cm,font=\footnotesize\sffamily,text=ink!70] at (-4,-7.35) {No one-shot or quorum. Primary is the fallback if the parse reading is unavailable.};
\node[anchor=north west,text width=7cm,align=left,font=\small\sffamily\bfseries] at (.45,-3.25) {SPECIALIZED LAYOUT\\SELECTED FIRST};
\pic[scale=.8] at (.7,-4.4) {dqlayout};
\node[align=left,text width=5cm,font=\footnotesize\sffamily] at (5.45,-5.05) {Equal-width columns or a routed ruled scan. This decision has priority over the general text/scan routes.};
\node[block] (d1) at (1.5,-6.45) {Column /\\ruled reader};
\node[block] (d2) at (4,-6.45) {Layout\\checks};
\node[block] (d3) at (6.5,-6.45) {Return its\\JSON};
\draw[flow] (d1)--(d2);\draw[flow] (d2)--(d3);
\node[align=left,text width=7cm,font=\footnotesize\sffamily,text=ink!70] at (4,-7.35) {One geometry-aware candidate. The API does not compare it with general-reader candidates.};
\end{tikzpicture}
